% ============================================================
% CHAPTER 2: SYSTEM REQUIREMENTS AND TECHNOLOGIES
% ============================================================

\chapter{System Requirements and Technologies}

Trustless Notes is a zero-knowledge encrypted note-taking application that combines modern web frameworks with client-side cryptography to ensure user data privacy. The system's requirements and technology choices are driven by the core architectural constraint: the server must never have access to plaintext data.

% ------------------------------------------------------------
\section{Software Requirements}
% ------------------------------------------------------------

\begin{itemize}
    \item \textbf{Frontend Framework:} Next.js 16 with React 19 and TypeScript provides the application shell, server-side rendering for landing pages, and API route handling. The App Router structure enables colocation of page components with their corresponding API logic.

    \item \textbf{User Interface:} Radix UI Themes with Tailwind CSS v4 provides the component library and styling system. The dark theme with red accent and slate neutral colors creates a security-focused aesthetic.

    \item \textbf{State Management:} Zustand manages client-side application state across three in-memory stores: session state (username and derived CryptoKey), notes list, and shared notes list. States are intentionally non-persistent to clear sensitive data on tab close.

    \item \textbf{Cryptography:} The Web Cryptography API (window.crypto.subtle) provides all cryptographic primitives: PBKDF2 for key derivation, AES-256-GCM for symmetric encryption, ECDH P-256 for key exchange, HKDF for shared secret derivation, and HMAC-SHA256 for cookie signing.

    \item \textbf{Database:} Supabase provides PostgreSQL-based relational storage for user accounts, encrypted note metadata, sharing records, and attachment metadata. Supabase Storage handles encrypted file blob storage.

    \item \textbf{Markdown Rendering:} react-markdown renders note content with GitHub-flavored Markdown support for rich text display.

    \item \textbf{Development Tools:} ESLint with eslint-config-next for code quality, TypeScript for type safety, and Git for version control.

    \item \textbf{Additional Libraries:} uuid for generating unique identifiers, cookie for server-side cookie parsing.
\end{itemize}

% ------------------------------------------------------------
\section{Hardware Requirements}
% ------------------------------------------------------------

The minimum hardware requirements for implementing and running the system are as follows:

\begin{itemize}
    \item A standard computer or laptop with a 64-bit operating system (Windows, macOS, or Linux).
    \item Minimum 8 GB RAM for running the development environment, Next.js development server, and testing.
    \item At least 2 GB of available storage for project files, dependencies (node\_modules), and build artifacts.
    \item Multi-core processor for efficient compilation and hot-reload during development.
    \item Stable internet connectivity for accessing Supabase cloud database, downloading dependencies, and testing real-time features.
    \item A modern web browser supporting the Web Cryptography API (Chrome, Firefox, Edge, or Safari).
\end{itemize}

% ------------------------------------------------------------
\section{Functional Requirements}
% ------------------------------------------------------------

The system serves registered users who can create, manage, and share encrypted notes. The following functional requirements define the system's behavior.

% ............................................................
\subsection{User Authentication}
% ............................................................

\begin{itemize}
    \item \textbf{Account Creation:}
    \begin{itemize}
        \item The system shall allow a user to create an account with a username and password.
        \item The system shall generate a random 16-byte salt on the client for PBKDF2 key derivation.
        \item The system shall derive a non-extractable AES-256-GCM key from the password using PBKDF2 with 200,000 iterations of SHA-256.
        \item The system shall generate a random sentinel value, encrypt it with the derived key, compute its SHA-256 hash, and store only the ciphertext, IV, and hash on the server.
        \item The system shall generate an ECDH P-256 keypair on the client, encrypt the private key with the derived key, and store the public key in plaintext and the private key as ciphertext on the server.
        \item The system shall send only cryptographic artifacts to the server; the password shall never be transmitted.
    \end{itemize}

    \item \textbf{Sign In:}
    \begin{itemize}
        \item The system shall always return an HTTP 200 response with indistinguishable sentinel data (real or fake) regardless of whether the username exists, preventing username enumeration.
        \item The system shall attempt sentinel decryption on the client using the locally derived key; if decryption fails, it shall display a generic "Invalid username or password" error.
        \item On successful decryption, the client shall send the decrypted sentinel plaintext to the server, which shall hash and compare it against the stored sentinel\_hash.
        \item The system shall issue an HMAC-SHA256 signed cookie containing the username upon successful sentinel verification.
    \end{itemize}

    \item \textbf{Session Management:}
    \begin{itemize}
        \item The system shall maintain session state via a signed `username` cookie.
        \item The system shall verify the cookie signature on every authenticated API request.
        \item The system shall allow session termination by clearing the cookie.
    \end{itemize}
\end{itemize}

% ............................................................
\subsection{Note Management}
% ............................................................

\begin{itemize}
    \item \textbf{Create Note:}
    \begin{itemize}
        \item The system shall generate a unique AES-256-GCM key for each new note.
        \item The system shall encrypt the note title and content separately using the note key.
        \item The system shall wrap (encrypt) the note key using the user's derived key.
        \item The system shall store the ciphertexts, IVs, and wrapped key on the server.
    \end{itemize}

    \item \textbf{Read Notes:}
    \begin{itemize}
        \item The system shall retrieve all encrypted notes belonging to the authenticated user.
        \item The client shall unwrap each note's key using the derived key and decrypt the title and content.
        \item The system shall display decrypted notes in a sidebar with search and selection.
    \end{itemize}

    \item \textbf{Update Note:}
    \begin{itemize}
        \item The system shall support updating the title and/or content of an existing note.
        \item The client shall re-encrypt only the changed fields with the existing note key.
        \item The system shall implement debounced auto-save with an 800ms delay.
        \item The server shall verify note ownership before allowing updates.
    \end{itemize}

    \item \textbf{Delete Note:}
    \begin{itemize}
        \item The system shall allow the note owner to delete a note and its associated attachments.
        \item The server shall verify note ownership before allowing deletion.
    \end{itemize}
\end{itemize}

% ............................................................
\subsection{Note Sharing}
% ............................................................

\begin{itemize}
    \item \textbf{Share Note:}
    \begin{itemize}
        \item The system shall allow the note owner to share a note with another registered user.
        \item The system shall fetch the recipient's ECDH public key (stored in plaintext on the server).
        \item The client shall decrypt its own ECDH private key using the derived key.
        \item The client shall compute the ECDH shared secret and derive an AES key via HKDF.
        \item The client shall wrap the note key with this derived AES key and store the ciphertext on the server.
        \item The server shall record the sharing relationship (sender, receiver, note, wrapped key).
    \end{itemize}

    \item \textbf{Receive Shared Note:}
    \begin{itemize}
        \item The system shall list all notes shared with the authenticated user.
        \item The client shall decrypt its own ECDH private key, compute the shared secret with the sender's public key, derive the AES key via HKDF, and unwrap the note key.
        \item Shared notes shall be displayed in a read-only editor.
    \end{itemize}
\end{itemize}

% ............................................................
\subsection{Attachment Management}
% ............................................................

\begin{itemize}
    \item \textbf{Upload Attachment:}
    \begin{itemize}
        \item The system shall allow users to attach image files to notes (maximum 5 MB per file).
        \item The client shall read the file as an ArrayBuffer, encrypt it with the note key using AES-256-GCM, and upload the encrypted blob to Supabase Storage.
        \item The system shall record metadata (IV, original filename, MIME type, storage path) in the attachments table.
    \end{itemize}

    \item \textbf{View and Download Attachments:}
    \begin{itemize}
        \item The system shall display thumbnail previews of attached images.
        \item On download, the client shall fetch the encrypted blob, decrypt it with the note key, and create an in-memory Blob URL for display.
        \item Decrypted data shall remain in memory only and be discarded when the component unmounts.
    \end{itemize}

    \item \textbf{Delete Attachment:}
    \begin{itemize}
        \item The system shall allow the note owner to delete individual attachments.
        \item The system shall remove both the database record and the Storage blob.
    \end{itemize}
\end{itemize}

% ............................................................
\subsection{Security Simulator}
% ............................................................

\begin{itemize}
    \item The system shall provide an interactive `/simulator` route with four attack scenarios.
    \item Each scenario shall visually demonstrate a specific breach and explain why it fails to compromise user data.
    \item Scenarios shall include: Database Leak, Password Leak, Cookie Forgery, and Supabase Service Key Leak.
    \item The simulator shall reimplement cryptographic operations with extractable keys for display purposes only.
    \item The simulator shall include a typewriter terminal component for narrative delivery.
\end{itemize}

% ------------------------------------------------------------
\section{Non-Functional Requirements}
% ------------------------------------------------------------

\begin{itemize}
    \item \textbf{Performance:}
    \begin{itemize}
        \item PBKDF2 key derivation (200,000 iterations) shall complete within 2 seconds on modern hardware.
        \item AES-256-GCM encryption and decryption of note content shall complete in under 100 milliseconds.
        \item Note list loading with decryption shall complete within 3 seconds for up to 100 notes.
        \item The debounced auto-save shall trigger no more than once per 800ms of inactivity.
    \end{itemize}

    \item \textbf{Security:}
    \begin{itemize}
        \item The user's derived key shall be non-extractable (cannot be exported from the Web Cryptography API).
        \item The user's password shall never be transmitted to the server.
        \item All API communication shall use HTTPS encryption.
        \item Session cookies shall be signed with HMAC-SHA256 to prevent forgery.
        \item The sign-in endpoint shall provide indistinguishability between valid and invalid usernames.
        \item Note ownership and sharing relationships shall be verified server-side on every operation.
        \item ECDH private keys shall be stored encrypted and never exposed in plaintext on the server.
    \end{itemize}

    \item \textbf{Reliability:}
    \begin{itemize}
        \item The system shall maintain data integrity through transaction management and foreign key constraints.
        \item Encrypted data shall be stored with authentication tags (AES-GCM provides authenticated encryption).
        \item The system shall gracefully handle network failures during saves with local state preservation.
    \end{itemize}

    \item \textbf{Usability:}
    \begin{itemize}
        \item The interface shall provide clear visual feedback during cryptographic operations (loading states).
        \item Error messages shall not reveal whether a username exists on the system.
        \item The markdown editor shall support standard formatting syntax.
    \end{itemize}
\end{itemize}
